% Generated draft: requires mathematical and source review.
\documentclass[11pt]{article}
\usepackage[margin=25mm]{geometry}
\usepackage{fontspec}
\setmainfont{Latin Modern Roman}
\usepackage{amsmath,amssymb,mathtools,unicode-math}
\setmathfont{Latin Modern Math}
\usepackage{enumitem,needspace,xurl,hyperref}
\setlength{\parindent}{0pt}
\setlength{\parskip}{5pt}
\setlength{\emergencystretch}{3em}
\newcommand{\N}{\mathbb{N}}
\newcommand{\Z}{\mathbb{Z}}
\newcommand{\Q}{\mathbb{Q}}
\newcommand{\R}{\mathbb{R}}
\newcommand{\C}{\mathbb{C}}
\newcommand{\F}{\mathbb{F}}
\newcommand{\A}{\mathbb{A}}
\newcommand{\PP}{\mathbb{P}}
\newcommand{\E}{\mathbb{E}}
\newcommand{\eps}{\varepsilon}
\newcommand{\dd}{\,\mathrm{d}}
\begin{document}
\section*{UnsolvedMath ID: 30002907}
\textbf{Source problem code:} catalogue entry
\subsection{Definitions and mathematical statement}
Let \(S_{g,n}\) be the surface denoted by this symbol in the supplied record, with \(g\geq 3\). Let \(K\subseteq\pi_1(S_{g,n})\) be a mapping-class-invariant subgroup of finite index defining a finite Galois covering
\[
S_K\longrightarrow S_{g,n}.
\]
Write
\[
H_K=H_1(\overline{S}_K;\Q),
\]
where \(\overline{S}_K\) denotes the compactified surface intended by the source. Let \(N_K\) denote the source's group of lifted mapping classes, and let
\[
\rho:N_K\longrightarrow\operatorname{GL}_{\Q}(H_K)
\]
be its action on homology. Write \(\mathcal G_K\) for the Zariski closure of \(\rho(N_K)\), and \(\mathcal G_K^\circ\) for its identity component. The original abbreviated statement uses \(G_K\) for this group.

The fixed subspace is
\[
H_K^{\mathcal G_K^\circ}
=\{v\in H_K: a\cdot v=v\text{ for every }a\in\mathcal G_K^\circ\},
\]
with invariance understood for the algebraic group action, after scalar extension when necessary. For a specified \(\Q\)-algebraic realization \(G\subseteq\operatorname{GL}(H_K)\), the usual meaning of an arithmetic subgroup is a subgroup commensurable with
\[
G(\Q)\cap\operatorname{GL}(L),
\]
where \(L\subset H_K\) is a full \(\Z\)-lattice and \(\operatorname{GL}(L)\) consists of the rational linear automorphisms preserving \(L\). Two subgroups are commensurable when their intersection has finite index in each.

The supplied questions ask whether, for every such covering with \(g\geq3\),
\[
H_K^{\mathcal G_K^\circ}=0,
\]
and whether \(\rho(N_K)\) is arithmetic in \(\mathcal G_K\).

These questions retain unresolved source conventions: the meaning and allowed values of \(n\), boundary and marking conventions for \(S_{g,n}\), the compactification \(\overline{S}_K\), and the precise definition of \(N_K\) are not supplied. Mapping-class invariance of \(K\) also requires a convention for the outer action on the fundamental group or a based formulation. The rational algebraic structure intended for \(\mathcal G_K\) is not stated explicitly; the arithmetic definition above explains the usual rational interpretation without resolving that omission. The questions therefore require these conventions to be recovered from the cited report before a fully unambiguous formulation can be accepted.
\subsection{Short English statement}
For every source-admissible finite Galois surface cover with base genus at least three, ask whether the connected algebraic monodromy group fixes no nonzero rational homology class and whether the lifted mapping classes have arithmetic monodromy image. Surface, lifting, compactification, and rational algebraic group conventions remain to be checked. This is a statement-only draft requiring human mathematical review.
\subsection{Sources}
\begin{itemize}
\item Local UnsolvedMath registry, permanent ID \(30002907\), title \emph{Invariant Prym Homology and Arithmetic Monodromy}, source problem code \texttt{catalogue entry}; collection: \emph{Oberwolfach Reports Open Problems}. The supplied record gives proposed year \(2015\) and names no proposer.
\item Supplied source citation: \emph{New Perspectives on the Interplay between Discrete Groups in Low-Dimensional Topology and Arithmetic Lattices} (2016), Oberwolfach Reports, pp.\ 1701--1746, DOI \texttt{10.4171/owr/2015/30}. Source URL and supplied literature-source link, labelled ``original Oberwolfach report'': \url{https://doi.org/10.4171/owr/2015/30}. The background describes Oberwolfach Reports as documenting workshops and their open-problem discussions. No exact problem passage or page within this report is supplied.
\item Supplied literature assessment dated \texttt{2026-08-21}: status ``open.'' The assessment describes unresolved questions about higher Prym homology and arithmetic image and says that the abbreviated notation does not support a stronger classification. This is supplied dated curation, not independently verified here; the background explicitly calls for independent verification of time-sensitive status and cited claims.
\item The supplied statement status is \texttt{corrected\_verified}. Its verification note reports that the official report passage was checked directly and that definitions and notation were restored without strengthening the question. This is a claim in the supplied record, not a source verification performed for this draft. The citation's year \(2016\), proposed year \(2015\), and DOI containing \texttt{2015} are retained as supplied.
\end{itemize}
\end{document}
